\documentclass[11pt]{article}

\usepackage[margin=0.95in]{geometry}
\usepackage{amsmath,amssymb,amsthm,mathtools,booktabs,enumitem,microtype}
\usepackage{xcolor}
\usepackage[colorlinks=true,linkcolor=blue!50!black,urlcolor=blue!50!black,citecolor=blue!50!black]{hyperref}

\theoremstyle{plain}
\newtheorem{theorem}{Theorem}
\newtheorem{lemma}[theorem]{Lemma}
\newtheorem{proposition}[theorem]{Proposition}
\newtheorem{corollary}[theorem]{Corollary}

\theoremstyle{definition}
\newtheorem{definition}[theorem]{Definition}

\theoremstyle{remark}
\newtheorem{remark}{Remark}

\newcommand{\Sow}{S_{\mathrm{ow}}}
\newcommand{\mkv}{\text{--}}

\title{\bfseries Logical Incompatibility of the Key-Agreement Conjectures\\and Exact Collapse of \texorpdfstring{$\Psi_\Delta$}{Psi\_Delta}}
\author{}
\date{\vspace{-1.6em}%
\small
\begin{tabular}{@{}c@{}}
Target paper: H.~Abin and A.~Gohari, \emph{On the Source Model Key Agreement Problem},\\
\emph{IEEE Trans.\ Inf.\ Theory} 72(8):5317--5334 (2026), \href{https://arxiv.org/abs/2502.00294}{arXiv:2502.00294}, Conjectures~\texttt{conjecture:1} and~\texttt{conj:2}\\[0.35em]
Companion verification script: \texttt{solutions/key\_agreement\_ineq/explore.py}\\[0.35em]
November 2026
\end{tabular}}

\begin{document}

\maketitle

\begin{abstract}
We resolve the joint relationship between Conjecture~1 (\texttt{conjecture:1}) and Conjecture~2 (\texttt{conj:2}) of Abin and Gohari (\emph{IEEE Trans.\ Inf.\ Theory}, 2026; \href{https://arxiv.org/abs/2502.00294}{arXiv:2502.00294}).
For every binary XOR source $(A, B, E)$ with $E = A \oplus B$, we construct an explicit six-variable common-observation instance of Conjecture~1 whose left-hand side minus right-hand side equals $I(A; B) - S(A; B \| A \oplus B)$ identically.
Consequently, Conjecture~1 implies the negation of Conjecture~2: the two conjectures cannot both hold, and any strict-gap witness to Conjecture~2 produces a counterexample to Conjecture~1 with the exact same gap.
More generally, Conjecture~1 forces $S(A; B \| E) = I(A; B)$ for every source satisfying the auxiliary-variable hypothesis of Theorem~\texttt{Th:gen:classic} in Abin and Gohari (2026).
Second, we evaluate the paper's constrained upper bound $\Psi_\Delta(A; B \| A \oplus B) = I(A; B)$ for every binary XOR law and every feasible $\Delta \in [0, H(E)]$ (alongside the Maurer--Wolf intrinsic information), refuting the paper's proposed strategy for proving Conjecture~2.
Neither conjecture is individually proved or refuted here.
We also characterize one-way equality, exhibit a strictly positive two-sample interactive rate for every dependent binary XOR source, and construct a triple-grouping protocol for the symmetric triangle source achieving $(1 + \log_2 3)/12 \approx 0.215414$ when both one-way capacities vanish.
\end{abstract}

\section{Model and Conjecture Statements}

All random variables take values in finite alphabets and all logarithms are base~$2$.
Alice, Bob, and Eve observe i.i.d.\ copies of a discrete memoryless source $(X, Y, Z)$, receiving $X^n, Y^n, Z^n$ respectively.
Using independent private randomness and an authenticated public interactive transcript $F$, Alice and Bob output keys $K_A, K_B$ in a common finite alphabet $\mathcal{K}$.
Following Ahlswede--Csisz\'ar~\cite{AhlswedeCsiszar1993} and Abin--Gohari~\cite{AbinGohari2026}, a rate $R \ge 0$ is achievable if for every $\delta > 0$ and sufficiently large $n$ there exists a protocol satisfying $\Pr[K_A \neq K_B] \le \delta$, $R = H(K_A)/n$,
\[
\frac{1}{n}\log_2 |\mathcal{K}| - R \le \delta, \qquad \frac{1}{n}I(K_A; Z^n, F) \le \delta.
\]
The secret-key capacity $S(X; Y \| Z)$ is the supremum of achievable rates.
The one-way capacity from $X$ to $Y$ is characterized by the single-letter formula~\cite{AhlswedeCsiszar1993,AbinGohari2026}
\[
\Sow(X \to Y \| Z) = \max_{(U, V) \mkv X \mkv (Y, Z)} \bigl[I(U; Y \mid V) - I(U; Z \mid V)\bigr],
\]
where standard cardinality bounds ensure the maximum is attained.
We use the operational inequalities $\Sow(X \to Y \| Z) \le S(X; Y \| Z) \le \min\{I(X; Y),\, I(X; Y \mid Z)\}$; in particular, $S(T; T \| E) = H(T \mid E)$, since setting $U = T$ and $V$ constant matches the upper bound $I(T; T \mid E) = H(T \mid E)$.

\begin{quote}
\textbf{Conjecture~1 (\texttt{conjecture:1}, \cite{AbinGohari2026}).}
\emph{For every finite joint distribution $P_{X, Y, Z, X', Y', Z'}$,}
\begin{equation}\label{eq:conj1}
\begin{split}
S(X; Y \| Z) - S(X'; Y' \| Z') \;\le\;{}& \Sow(X, Y \to Z' \| Z) + I(Y', Z'; X \mid X') \\
&+ I(X', Z'; Y \mid Y') + I(X; Y \mid X', Y', Z').
\end{split}
\end{equation}
\end{quote}

\begin{quote}
\textbf{Conjecture~2 (\texttt{conj:2}, \cite{AbinGohari2026}).}
\emph{There exists a binary joint distribution $P_{X, Y}$ on $\{0, 1\}^2$ such that $S(X; Y \| X \oplus Y) < I(X; Y)$.}
\end{quote}

\section{Joint Incompatibility: Conjecture~1 Implies \texorpdfstring{$\neg$}{not }Conjecture~2}

\begin{lemma}[Common-observation identity and XOR lift]\label{lem:lift}
\leavevmode
\begin{enumerate}[label=\textup{(\arabic*)}]
\item For any finite joint distribution $P_{T, A, B, E}$,
\begin{equation}\label{eq:identity}
\begin{split}
H(T \mid E) - H(T \mid A) - H(T \mid B) + H(T \mid A, B, E)
&= I(A; B) - I(A; B \mid T) \\
&\quad - I(T; E) - I(T; E \mid A, B).
\end{split}
\end{equation}
\item For every binary distribution $P_{A, B}$ on $\{0, 1\}^2$ with $E = A \oplus B$, there exists a random variable $T$ on a $4$-element alphabet coupled to $(A, B)$ such that
\begin{equation}\label{eq:lift-props}
I(A; B \mid T) = 0, \qquad I(T; E) = 0, \qquad I(T; E \mid A, B) = 0, \qquad H(A, B \mid T, E) = 0.
\end{equation}
\end{enumerate}
\end{lemma}

\begin{proof}
For~(1), expanding every conditional entropy and mutual information into joint entropies gives
\begin{align*}
\text{LHS} &= H(T, E) - H(E) - H(T, A) + H(A) - H(T, B) + H(B) + H(T, A, B, E) - H(A, B, E), \\
\text{RHS} &= \bigl[H(A) + H(B) - H(A, B)\bigr] - \bigl[H(T, A) + H(T, B) - H(T, A, B) - H(T)\bigr] \\
&\quad - \bigl[H(T) + H(E) - H(T, E)\bigr] - \bigl[H(T, A, B) + H(A, B, E) - H(T, A, B, E) - H(A, B)\bigr],
\end{align*}
which agree identically.

For~(2), denote the probability masses of $(A, B)$ at $(00, 01, 10, 11)$ by $(a, b, c, d)$, set $s := a + d$ and $t := b + c$, and define $u := a/s$ (setting $u := 0$ if $s = 0$) and $v := b/t$ (setting $v := 0$ if $t = 0$).
Define $T \in \{1, 2, 3, 4\}$ with weights and conditional laws for $(A, B)$ listed in Table~\ref{tab:lift}.
The weighted column sums equal $(su,\, tv,\, t(1 - v),\, s(1 - u)) = (a, b, c, d)$, including the boundary cases $s = 0$ or $t = 0$.
In every state of $T$, one coordinate of $(A, B)$ is constant, so $I(A; B \mid T) = 0$; the parity $E = A \oplus B$ has distribution $(s, t)$ in every state, so $I(T; E) = 0$; $E$ is a function of $(A, B)$, so $I(T; E \mid A, B) = 0$; and knowing $T$ and $E$ fixes one bit and the XOR, hence determines $(A, B)$, giving $H(A, B \mid T, E) = 0$.
\end{proof}

\begin{table}[h]
\centering
\begin{tabular}{@{}cllc@{}}
\toprule
State $T$ & Weight $P_T$ & Conditional law on $(00, 01, 10, 11)$ & Fixed bit \\
\midrule
$1$ & $u v$ & $(s,\, t,\, 0,\, 0)$ & $A = 0$ \\
$2$ & $u(1 - v)$ & $(s,\, 0,\, t,\, 0)$ & $B = 0$ \\
$3$ & $(1 - u) v$ & $(0,\, t,\, 0,\, s)$ & $B = 1$ \\
$4$ & $(1 - u)(1 - v)$ & $(0,\, 0,\, t,\, s)$ & $A = 1$ \\
\bottomrule
\end{tabular}
\caption{Product-mixture lift $P_{T, A, B}$ for a binary law $(a, b, c, d)$ with $s = a + d$, $t = b + c$.}\label{tab:lift}
\end{table}

\begin{theorem}[Incompatibility of Conjectures~1 and~2]\label{thm:incomp}
Assume Conjecture~1 holds.
Then every finite source $(A, B, E)$ admitting an auxiliary variable $T$ with $I(A; B \mid T) = I(T; E) = I(T; E \mid A, B) = 0$ (i.e., every source satisfying Theorem~\texttt{Th:gen:classic} of~\cite{AbinGohari2026}) satisfies $S(A; B \| E) = I(A; B)$.
In particular, Conjecture~1 implies $S(A; B \| A \oplus B) = I(A; B)$ for every binary law $P_{A, B}$, which is the negation of Conjecture~2.
\end{theorem}

\begin{proof}
Apply~\eqref{eq:conj1} with $(X, Y, Z) = (T, T, E)$ and $(X', Y', Z') = (A, B, E)$.
Then $S(T; T \| E) = H(T \mid E)$, and $\Sow(T, T \to E \| E) = 0$ because $I(U; E \mid V) - I(U; E \mid V) = 0$ identically.
The remaining three penalty terms evaluate to
\begin{align*}
&I(B, E; T \mid A) + I(A, E; T \mid B) + I(T; T \mid A, B, E) \\
&\quad = \bigl[H(T \mid A) - H(T \mid A, B, E)\bigr]
  + \bigl[H(T \mid B) - H(T \mid A, B, E)\bigr] + H(T \mid A, B, E) \\
&\quad = H(T \mid A) + H(T \mid B) - H(T \mid A, B, E).
\end{align*}
Subtracting the right-hand side of~\eqref{eq:conj1} from its left-hand side and invoking~\eqref{eq:identity} yields
\[
\text{LHS}\eqref{eq:conj1} - \text{RHS}\eqref{eq:conj1}
= I(A; B) - S(A; B \| E) - I(A; B \mid T) - I(T; E) - I(T; E \mid A, B).
\]
When the last three terms vanish, \eqref{eq:conj1} forces $S(A; B \| E) \ge I(A; B)$; combining with the standard upper bound $S(A; B \| E) \le I(A; B)$ gives $S(A; B \| E) = I(A; B)$.
Applying Lemma~\ref{lem:lift}(2) completes the proof for all binary XOR sources.
\end{proof}

\begin{remark}[Scope caveat and restricted equivalence]
Neither Conjecture~1 nor Conjecture~2 is individually settled here.
Theorem~\ref{thm:incomp} proves that Conjecture~1 holds on the restricted family of common-observation XOR instances $((T, T, E), (A, B, E))$ if and only if Conjecture~2 is false: any witness with $S(A; B \| A \oplus B) < I(A; B)$ immediately refutes Conjecture~1 with the identical gap, whereas disproving Conjecture~2 would still leave arbitrary six-variable instances of Conjecture~1 open.
\end{remark}

\section{Exact Collapse of \texorpdfstring{$\Psi_\Delta$}{Psi\_Delta} and Intrinsic Information}

To pursue Conjecture~2, Abin and Gohari~\cite{AbinGohari2026} proposed the constrained upper bound $\Psi_\Delta(A; B \| E)$ for protocols whose normalized residual XOR entropy is at least $\Delta \in [0, H(E)]$:
\[
\Psi_\Delta(A; B \| E) := \inf_{P_{J \mid A, B, E}} \max_{\substack{(U, V) \mkv (A, B) \mkv (J, E) \\ I(V; E) \le I(V; J) + H(E) - \Delta}} \bigl[I(A; B \mid J) + I(U; J \mid V) - I(U; E \mid V)\bigr],
\]
allowing the additional XOR constraint $H(A, B \mid U, V, E) = 0$ and binary cardinality bound $|\mathcal{U}| \le 4$.

\begin{theorem}[Collapse of $\Psi_\Delta$ and intrinsic information]\label{thm:psi}
For every binary law $P_{A, B}$ with $E = A \oplus B$ and every $0 \le \Delta \le H(E)$,
\begin{equation}\label{eq:psi}
\Psi_\Delta(A; B \| E) = I(A; B).
\end{equation}
Moreover, the Maurer--Wolf intrinsic information $I(A; B \downarrow E) := \inf_{(A, B) \mkv E \mkv J} I(A; B \mid J)$ equals $I(A; B)$.
\end{theorem}

\begin{proof}
Choosing $J$ constant yields $\Psi_\Delta(A; B \| E) \le I(A; B)$ since $I(U; J \mid V) = 0$ and $I(U; E \mid V) \ge 0$.
Conversely, fix any finite channel $P_{J \mid A, B}$ (equivalent to $P_{J \mid A, B, E}$ since $E = A \oplus B$).
Generate the $4$-state variable $T$ of Lemma~\ref{lem:lift}(2) from its conditional law given $(A, B)$, conditionally independent of $J$ given $(A, B)$, so that $I(T; J \mid A, B) = 0$.
Set $U = T$ and take $V$ constant.
Then $(U, V) \mkv (A, B) \mkv (J, E)$, $|\mathcal{U}| \le 4$, $|\mathcal{V}| = 1$, $H(A, B \mid U, V, E) = H(A, B \mid T, E) = 0$, and $I(V; E) = 0 \le I(V; J) + H(E) - \Delta$ throughout $[0, H(E)]$.

Since $I(A; B \mid T) = 0$,
\[
I(J; A \mid T) + I(J; B \mid T) = I(J; A, B \mid T) - I(A; B \mid J, T) \le I(J; A, B \mid T).
\]
Applying the chain rule with $I(T; J \mid A, B) = 0$ therefore gives
\begin{align*}
I(A; B) - I(A; B \mid J)
&= I(J; A) + I(J; B) - I(J; A, B) \\
&= \bigl[I(J; T) + I(J; A \mid T) - I(J; T \mid A)\bigr] + \bigl[I(J; T) + I(J; B \mid T) - I(J; T \mid B)\bigr] \\
&\quad - \bigl[I(J; T) + I(J; A, B \mid T) - I(J; T \mid A, B)\bigr] \\
&\le I(J; T) - I(J; T \mid A) - I(J; T \mid B) \le I(J; T).
\end{align*}
Since $I(T; E) = 0$, this feasible $(U, V)$ achieves $I(A; B \mid J) + I(T; J) - I(T; E) \ge I(A; B)$ for every channel $P_{J \mid A, B}$, proving~\eqref{eq:psi}.
Finally, for any degradation $(A, B) \mkv E \mkv J$, the chain $T \mkv E \mkv J$ and $I(T; E) = 0$ force $I(J; T) = 0$, so the same inequality gives $I(A; B \mid J) \ge I(A; B)$, attained by constant~$J$.
\end{proof}

Theorem~\ref{thm:psi} evaluates the upper bound $\Psi_\Delta$ identically to $I(A; B)$ across all $\Delta \in [0, H(E)]$, even if the operational subclass at a given $\Delta$ had strictly smaller capacity.
Thus neither optimizing Eve's degradation~\cite{MaurerWolf1999} nor strengthening the residual-entropy parameter $\Delta$ in $\Psi_\Delta$~\cite{GohariAnantharam2010,AbinGohari2026} can certify Conjecture~2.

\section{Operational Rates for Binary XOR Sources}

\begin{proposition}[One-way equality criterion]\label{prop:ow}
Let $(A, B)$ be dependent binary random variables ($I(A; B) > 0$) and $E = A \oplus B$.
Then $\Sow(A \to B \| E) = I(A; B)$ if and only if $A \perp E$.
By symmetry, $\Sow(B \to A \| E) = I(A; B)$ if and only if $B \perp E$.
\end{proposition}

\begin{proof}
For any $(U, V)$ satisfying $(U, V) \mkv A \mkv (B, E)$, the chain rule gives the exact gap
\[
I(A; B) - \bigl[I(U; B \mid V) - I(U; E \mid V)\bigr] = I(V; B) + I(A; B \mid U, V) + I(U; E \mid V).
\]
Because $A, B$ are binary and dependent, the conditional distributions $P_{B \mid A = 0}$ and $P_{B \mid A = 1}$ are distinct.
Thus $I(V; B) = 0$ uniquely fixes $P(A = 1 \mid V) = P(A = 1)$, i.e., $V \perp A$.
Likewise, $I(A; B \mid U, V) = 0$ forces $A$ to be determined by $(U, V)$, since otherwise some $(u, v)$ supports both distinct conditional laws of $B$.
Then $(U, V) \mkv A \mkv E$ and $V \perp A$ yield $I(U; E \mid V) = I(A; E \mid V) = I(A; E)$.
Since the one-way maximum is attained, equality forces $I(A; E) = 0$.
Conversely, if $A \perp E$, setting $U = A$ and $V$ constant attains $I(A; B)$.
\end{proof}

\begin{proposition}[Strictly positive two-sample interactive rate]\label{prop:twosample}
For every binary XOR source with masses $(a, b, c, d)$ at $(00, 01, 10, 11)$,
\[
S(A; B \| A \oplus B) \;\ge\; (ad + bc)\left[1 - h_2\!\left(\frac{ad}{ad + bc}\right)\right]
\]
(interpreted as $0$ when $ad + bc = 0$), which is strictly positive if and only if $ad \neq bc$.
\end{proposition}

\begin{proof}
Group $2m$ i.i.d.\ samples into $m$ pairs.
Alice announces $A_{2i-1} \oplus A_{2i}$ and Bob announces $B_{2i-1} \oplus B_{2i}$ for each pair $i$, retaining pairs where both parities equal~$1$ (probability $2(ad + bc)$ per pair).
Given acceptance, $(A_{2i-1}, B_{2i-1})$ has law proportional to $(ad, bc, bc, ad)$ on $(00, 01, 10, 11)$: both marginal bits are fair, and Eve's two observations $(E_{2i-1}, E_{2i})$ are identical and independent of $A_{2i-1}$.
Retaining a deterministic count $\lfloor m \cdot 2(ad + bc)(1 - \varepsilon)\rfloor$ of accepted pairs (aborting with exponentially small probability on shortage) and running one-way key extraction from $A$ to $B$ yields $1 - h_2(ad/(ad + bc))$ bits per accepted pair, hence $(ad + bc)[1 - h_2(ad/(ad + bc))]$ per original sample.
\end{proof}

\begin{proposition}[Triangle source]\label{prop:triangle}
For the binary XOR source $(a, b, c, d) = (1/3, 1/3, 1/3, 0)$, both one-way capacities vanish, $\Sow(A \to B \| E) = \Sow(B \to A \| E) = 0$, whereas
\[
\frac{1 + \log_2 3}{12} = 0.215413541727\ldots \;\le\; S(A; B \| E) \;\le\; I(A; B) = \log_2 3 - \frac{4}{3} = 0.251629167388\ldots
\]
Moreover, the common-observation instance of Theorem~\ref{thm:incomp} uses a single fair bit $T$ (all six variables binary), with unprimed capacity~$1$ and total penalty $7/3 - \log_2 3$.
\end{proposition}

\begin{proof}
Given $A$, the conditional laws of $B$ and $1 - E$ coincide, so $I(U; B \mid V) = I(U; E \mid V)$ for all $(U, V) \mkv A \mkv (B, E)$, giving $\Sow(A \to B \| E) = 0$; symmetry gives $\Sow(B \to A \| E) = 0$.
Setting $T = 0$ at $(0, 1)$, $T = 1$ at $(1, 0)$, and a fair coin at $(0, 0)$ yields a binary $T \sim \mathrm{Bern}(1/2)$ with $I(A; B \mid T) = I(T; E) = I(T; E \mid A, B) = 0$ and $H(T \mid E) = 1$.

For the interactive lower bound, Alice independently shuffles her lists of $1$-positions and $0$-positions in $A^n$, partitions as many positions as possible into $m = \min(N_1, \lfloor N_0 / 2 \rfloor)$ disjoint triples each containing one~$1$ and two~$0$s ($m/n \to 1/3$ in probability), and publishes $N_1$ together with the unordered triples.
In each triple, call the index of Alice's~$1$ the \emph{pointer}.
Permuting the three coordinates inside any triple preserves both the source probability and the grouping likelihood, so independent within-triple permutations act transitively on pointer tuples conditional on $N_1$, the unordered triples, and the ungrouped symbols; hence the pointers are i.i.d.\ uniform trits.

In each triple, Bob holds $0$ at the pointer and two independent fair bits elsewhere, and announces his Hamming weight:
\begin{itemize}[leftmargin=1.5em]
\item \emph{Weight~$2$ (probability $1/4$):} Bob identifies the pointer uniquely and Eve sees $(1, 1, 1)$, yielding a uniform secret trit independent of Eve's view and the transcript.
\item \emph{Weight~$1$ (probability $1/2$):} Alice's pointer, Bob's~$1$, and Eve's~$0$ occupy the three positions in a uniform random order. Alice publishes an unordered pair consisting of her pointer and one uniformly chosen other position; Bob announces success if his~$1$ belongs to the pair (probability $1/2$, total mass $1/4$). On success, both know the pointer, while Eve's~$0$ lies outside the pair and both endpoints remain equiprobable given Eve's view and the transcript, yielding one uniform secret bit.
\item \emph{Remaining cases (probability $1/2$):} Discarded.
\end{itemize}
Retaining deterministic counts $M_1 = \lfloor m(1/4 - \varepsilon) \rfloor$ of weight-$2$ triples and $M_2 = \lfloor m(1/4 - \varepsilon) \rfloor$ of successful weight-$1$ triples gives vanishing abort probability and a uniform secret key in $\mathcal{K} = \{0, 1, 2\}^{M_1} \times \{0, 1\}^{M_2}$ independent of $(E^n, F)$.
Under the finite-alphabet definition of~\cite{AbinGohari2026} ($\frac{1}{n}\log_2 |\mathcal{K}| - R \le \delta$), outputting these $M_1$ trits and $M_2$ bits directly achieves rate $\frac{1}{3}(\frac{1}{4}\log_2 3 + \frac{1}{4}) = (1 + \log_2 3)/12$ without needing trit-to-bit binning (and optional binning to binary strings incurs only $o(1)$ rate loss).
\end{proof}

\section{Reproducible Verification}

The script \texttt{solutions/key\_agreement\_ineq/explore.py} enumerates all $455$ binary laws with denominator~$12$, verifies the $4$-state lift and Theorem~\ref{thm:incomp}'s gap identity on every law, checks $36{,}855$ binary degradation channels and $2{,}275$ random $3$-output auxiliary channels for Theorem~\ref{thm:psi}, confirms that all $390$ dependent grid laws have a strictly positive two-sample rate (Proposition~\ref{prop:twosample}), checks $500$ seeded six-variable Conjecture~1 instances with analytically known capacities, and verifies the exact branch distributions for Proposition~\ref{prop:triangle}. All assertions pass.

\begin{thebibliography}{9}

\bibitem{AbinGohari2026}
H.~Abin and A.~Gohari.
On the source model key agreement problem.
\emph{IEEE Transactions on Information Theory}, 72(8):5317--5334, August 2026.
Preprint \href{https://arxiv.org/abs/2502.00294}{arXiv:2502.00294}.

\bibitem{AhlswedeCsiszar1993}
R.~Ahlswede and I.~Csisz\'ar.
Common randomness in information theory and cryptography---Part~I: Secret sharing.
\emph{IEEE Transactions on Information Theory}, 39(4):1121--1132, 1993.

\bibitem{GohariAnantharam2010}
A.~Gohari and V.~Anantharam.
Information-theoretic key agreement of multiple terminals---Part~I.
\emph{IEEE Transactions on Information Theory}, 56(8):3973--3996, 2010.

\bibitem{MaurerWolf1999}
U.~M.~Maurer and S.~Wolf.
Unconditionally secure key agreement and the intrinsic conditional information.
\emph{IEEE Transactions on Information Theory}, 45(2):499--514, 1999.

\end{thebibliography}

\end{document}
